\chapter{اللقاء العاشر: قصة آدم والملائكة}

\section{مراجعة لقصة استخلاف آدم وسؤال الملائكة}
السلام عليكم ورحمة الله وبركاته. الحمد لله رب العالمين، والصلاة والسلام على نبينا محمد وعلى آله وصحبه أجمعين.

نبدأ بمراجعة لخلاصة كلام الطبري في قوله تعالى: \begin{ayah}إِنِّي جَاعِلٌ فِي الْأَرْضِ خَلِيفَةً ۖ قَالُوا أَتَجْعَلُ فِيهَا مَن يُفْسِدُ فِيهَا وَيَسْفِكُ الدِّمَاءَ\end{ayah}.
وقد لخصنا سابقاً أن الطبري رجح أن سؤال الملائكة كان على سبيل "الاستخبار والاستعلام" وليس على وجه الإنكار. واستدل على أن الملائكة علمت بوقوع الفساد إما بإعلام الله لهم، أو قياساً على من سكن الأرض قبلهم من الجن.

\separator

\section{تأويل (ونحن نسبح بحمدك ونقدس لك)}
\isnad{قال الإمام الطبري رحمه الله:} «أما قوله \textit{ونحن نسبح بحمدك} فإنه يعني: إنا نعظمك بالحمد لك والشكر... وكل ذكر لله عند العرب فتسبيح وصلاة... وأصل التسبيح لله عند العرب: التنزيه له من إضافة ما ليس من صفاته إليه، والتبرئة له من ذلك».

\begin{fawaid}[دلالة لفظ التسبيح]
لفظ التسبيح لا ينحصر في "مجرد التنزيه" ونفي النقص عن الله (وإن كان هذا من معانيه الكثيرة المستعملة)، ولكنه يأتي كذلك بمعنى التعظيم والذكر والصلاة، كما في قوله تعالى: \begin{ayah}سُبْحَانَ الَّذِي أَسْرَىٰ بِعَبْدِهِ لَيْلًا\end{ayah}.
\end{fawaid}

\isnad{وأما قوله:} \begin{ayah}وَنُقَدِّسُ لَكَ\end{ayah} \isnad{قال الطبري:} «والتقديس هو التطهير والتعظيم، ومنه قولهم: سُبُّوح قُدُّوس... ومعنى قول الملائكة: ننزهك ونبرئك مما يضيفه إليك أهل الشرك، ونصلي لك ونطهرك».

\separator

\section{تأويل (إني أعلم ما لا تعلمون)}
ساق \rawi{الطبري} الأقوال في مراد الله بذلك، فذكر أن بعضهم قال: يعلم من إبليس المعصية والكبر وإضماره لخلاف طاعة الله، وقال آخرون: أعلم أنه سيكون من ذرية هذا الخليفة أنبياء ورسل وقوم صالحون.

\begin{fawaid}[تحرير محل الخلاف في المسألة]
معلوم ومُتفق عليه أن الله يعلم كل غيب ولا نعلمه، ولكن المفسرين حينما يختلفون هنا فإنهم يبحثون في تحديد "هذا المعنى الغائب المعين" الذي قصدته الآية في هذا السياق تحديداً.
\end{fawaid}

\separator

\section{تأويل (وعلم آدم الأسماء كلها ثم عرضهم على الملائكة)}
اختلفوا في الأسماء التي علمها آدم، هل هي أسماء كل شيء (كالجبال والبحار والحيوانات)، أم أسماء الملائكة، أم أسماء ذريته؟

\isnad{قال أبو جعفر:} «وأولى هذه الأقوال بالصواب... قول من قال: إنها أسماء ذريته وأسماء الملائكة دون أسماء سائر أجناس الخلق؛ وذلك أن الله تعالى ذكره قال: \textit{ثم عرضهم} يعني بذلك أعيان المسمين، ولا تكاد العرب تكنى بالهاء والميم (هم) إلا عن أسماء بني آدم والملائكة».

\begin{fawaid}[الاستشهاد بالقراءات الشاذة على صحة المعنى اللغوي]
أشار \rawi{الطبري} إلى قراءة \rawi{ابن مسعود} (ثم عرضهن) وقراءة \rawi{أبي بن كعب} (ثم عرضها)، واستأنس بهما لبيان أن من فسر الأسماء بأنها "أسماء كل شيء" كان متوجهاً من الناحية اللغوية بناءً على هذه القراءات التفسيرية، ولكنه رجح تخصيصها بالعقلاء بناءً على قراءة المصحف العثماني المجمع عليه (ثم عرضهم).
\end{fawaid}

\separator

\section{تأويل (فقال أنبئوني بأسماء هؤلاء إن كنتم صادقين)}
\isnad{قال أبو جعفر:} «وتأويل قوله \textit{أنبئوني}: أخبروني... \textit{إن كنتم صادقين}: إن كنتم صادقين أني إن جعلت خليفتي في الأرض من غيركم عصاني ذريته وأفسدوا فيها، وإن جعلتكم فيها أطعتموني... فإنكم إذ كنتم لا تعلمون أسماء هؤلاء الذين عرضتهم عليكم... فأنتم بما هو غير موجود أحْرى أن تكونوا غير عالمين».

\subsection{رد الطبري على تأويل اللغويين}
\isnad{قال الطبري:} «وقد زعم بعض نحويي أهل البصرة [يقصد أبا عبيدة معمر بن المثنى] أن قوله: \textit{إن كنتم صادقين} لم يكن ذلك لأن الملائكة ادعوا شيئاً، إنما أخبر عن جهلهم... كما يقول الرجل: أنبئني بهذا إن كنت تعلم وهو يعلم أنه لا يعلم، يريد أنه جاهل».

فرد عليه \rawi{الطبري} بقوة مبيناً أن الصدق يقابل الكذب، ويكون في الأخبار والأفعال، ولا يجوز لغةً أن يقال "صدق الرجل" بمعنى "علم". وخطّأه بخروجه عن إجماع المفسرين.

\begin{fawaid}[قاعدة في الرد على الأقوال]
إذا ذكر العالم المحقق قولاً ليرده، فإنه يشرح أولاً وجه هذا القول ومستنده، ثم يشرع في نقضه ببيان لوازمه الباطلة ومخالفته للغة العرب ولإجماع المفسرين (أهل التأويل).
\end{fawaid}

\separator

\section{تأويل إبليس واستكباره}
\isnad{قال جل ثناؤه:} \begin{ayah}وَإِذْ قُلْنَا لِلْمَلَائِكَةِ اسْجُدُوا لِآدَمَ فَسَجَدُوا إِلَّا إِبْلِيسَ أَبَىٰ وَاسْتَكْبَرَ وَكَانَ مِنَ الْكَافِرِينَ\end{ayah}.
اختلفوا في إبليس: هل هو من الملائكة أم من غيرهم؟
ساق \rawi{الطبري} الأقوال، فرأى قوم أنه من حي من الملائكة يقال لهم الجن خُلقوا من نار السموم (وهو مروي عن ابن عباس)، ورأى آخرون (كالحسن البصري) أنه لم يك من الملائكة طرفة عين بل هو أصل الجن.

\isnad{ورجح الطبري} القول الأول مستدلاً بالاستثناء المتصل (فسجدوا إلا إبليس)، وراداً على من زعم أن الملائكة لا تتناسل ولا تخلق من نار، بأنه لا يمتنع أن يخلق الله صنفاً من الملائكة من نار ويثبت لهم التناسل.

\begin{fawaid}[أصل كلمة إبليس وتصريفه]
إبليس (إفْعِيل) من الإبلاس، وهو الإياس من الخير والندم والحزن. ومُنع من الصرف (لم يُنَوّن) استثقالاً، إذ شُبه بأسماء العجم لعدم وجود نظير له في أسماء العرب، فجرى مجرى إبراهيم وإسحاق.
\end{fawaid}

\separator

\section{تأويل الشجرة التي نُهي عنها آدم}
\isnad{قال أبو جعفر:} «ثم اختلف أهل التأويل في عين الشجرة التي نُهي عن أكل ثمارها آدم... فقال بعضهم: هي السنبلة (البر)، وقال آخرون: هي الكرمة (العنب)، وقال آخرون: هي التينة».

\begin{fawaid}[منهج الطبري في تفاصيل القصص المبهمة]
أصّل \rawi{الطبري} قاعدة ذهبية في التعامل مع هذه التفاصيل المُبهمة في القرآن؛ فبعد سرد الأقوال قال: «والصواب في ذلك أن يقال: إن الله نهى آدم عن شجرة بعينها... ولا عِلم عندنا بأي ذلك من أي، وجائز أن تكون واحدة منها. وذلك عِلمٌ إذا عُلِم لم ينفع، وإن جُهِل لم يَضُر». فما سكت عنه الوحي ولم يبينه فلا فائدة ولا طائل في البحث عنه، والعبرة متحققة بدونه.
\end{fawaid}

\separator

\section{تأويل (فأزلهما الشيطان عنها)}
اختلفت القراءة في ذلك؛ فقرأها عامة القراء: \textit{فأَزَلَّهُمَا} (بتشديد اللام) من الزلل وهو الوقوع في الخطيئة، وقرأها آخرون: \textit{فأَزَالَهُمَا} (بالألف وتخفيف اللام) بمعنى التنحية والإخراج.

\begin{fawaid}[الترجيح بين القراءات بالمعنى التأسيسي]
رجح \rawi{الطبري} قراءة التشديد (فأزلهما) لأن الله قال بعدها: (فأخرجهما)، فلو قرأنا (فأزالهما) لكان المعنى مكرراً (نحّاهما فأخرجهما). والقاعدة أن القراءة التي تؤسس لمعنى جديد (استذلهما فأوقعهما في الخطيئة، وبسبب ذلك أخرجهما) أولى من القراءة التي تفيد التأكيد المجرد.
\end{fawaid}

\separator

\section{كيف وسوس الشيطان لآدم؟ (موقف الطبري من الإسرائيليات)}
أورد \rawi{الطبري} القصة الإسرائيلية الطويلة المشهورة عن إدخال الحية لإبليس في فمها وتسلله إلى الجنة، ثم عقب عليها بتأصيل منهجي عظيم.

\begin{fawaid}[منهج الطبري في نقد الإسرائيليات]
رد \rawi{الطبري} تفاصيل القصة المزعومة محتجاً بظاهر القرآن: \begin{ayah}وَقَاسَمَهُمَا إِنِّي لَكُمَا لَمِنَ النَّاصِحِينَ\end{ayah}، فالمقاسمة والمخاطبة تقتضي المباشرة. وقرر قاعدته: «غير جائز أن يُترك المفهوم من ظاهر الكتاب... إلى باطل لا دلالة عليه من ظاهر التنزيل، ولا خبر عن الرسول ﷺ، ولا فيه إجماع مستفيض». فالإسرائيليات تُرد وتُطرح إذا خالفت ظاهر القرآن.
\end{fawaid}

\separator

\section{تأويل (فتلقى آدم من ربه كلمات فتاب عليه)}
قُرئت (فَتَلَقَّى آدَمَ مِنْ رَبِّهِ كَلِمَاتٌ) بنصب آدم ورفع كلمات على أن الكلمات هي التي تلقته.

\isnad{قال الطبري مقيماً الحجة على ردها:} «وذلك وإن كان من جهة العربية جائزاً... فغير جائز عندي في القراءة إلا رفع آدم... لإجماع الحجة من القرأة وأهل التأويل... وغير جائز الاعتراض عليها بقول من يجوز عليه السهو والخطأ».

\begin{fawaid}[رد القراءة الشاذة رغم صحتها لغوياً]
لا يكفي لصحة القراءة مجرد صحة وجهها في العربية واتساقها مع القواعد النحوية، بل لابد من موافقتها لنقل الحجة المستفيض وإجماع القرأة وعلماء الأمة.
\end{fawaid}

واختلفوا في الكلمات التي تلقاها آدم، ورجح \rawi{الطبري} أنها قولهما: \begin{ayah}رَبَّنَا ظَلَمْنَا أَنفُسَنَا وَإِن لَّمْ تَغْفِرْ لَنَا وَتَرْحَمْنَا لَنَكُونَنَّ مِنَ الْخَاسِرِينَ\end{ayah}؛ لأن هذا هو الثابت والمنصوص عليه في كتاب الله بسورة الأعراف، وما سواه محتمل لا حجة قاطعة عليه.

\separator

بهذا ننهي هذا اللقاء من تفسير سورة البقرة، ونكمل في المجلس القادم بإذن الله.